% Luc Destombe --- licence CC BY-NC-SA 4.0
% https://creativecommons.org/licenses/by-nc-sa/4.0/deed.fr
% Fichier autonome : compiler avec pdflatex, deux fois (signets, nombre de pages).
\documentclass[a4paper,10pt]{article}
\makeatletter
\newif\iflycee@niveaux
\lycee@niveauxtrue
\newif\iflycee@palatino
\lycee@palatinotrue

\RequirePackage[utf8]{inputenc}
\RequirePackage[T1]{fontenc}
\RequirePackage[french]{babel}
\RequirePackage{etoolbox}

\newcommand{\ShorthandsOff}{\@ifundefined{shorthandoff}{}{\shorthandoff{;:!?}}}
\newcommand{\ShorthandsOn}{\@ifundefined{shorthandon}{}{\shorthandon{;:!?}}}
\ShorthandsOff
\AtBeginDocument{\ShorthandsOn}
\AtBeginEnvironment{tikzpicture}{\ShorthandsOff}

\RequirePackage{geometry}
\RequirePackage{fancyhdr}
\RequirePackage{lastpage}
\RequirePackage{multicol}
\RequirePackage{array}
\RequirePackage{enumitem}
\RequirePackage{xcolor}
\RequirePackage{needspace}

\RequirePackage{amsmath,amssymb}
\RequirePackage{mathpazo}
\DeclareMathSizes{10.95}{10.95}{8}{5.5}
\begingroup\catcode`\_=\active \catcode`\^=\active
\gdef_#1{\sb{\scriptscriptstyle #1}}
\gdef^#1{\sp{\scriptscriptstyle\lycee@moinsepais #1}}
\endgroup
\newcommand\lycee@moins{\mathbin{%
  \setbox\z@\hbox{$\m@th\scriptscriptstyle\lycee@moinsorig$}%
  \dimen@=.55\wd\z@ \advance\dimen@-.5pt
  \hbox to.75\wd\z@{\hss$\m@th\scriptscriptstyle\vcenter{\hbox{%
    \tikz\draw[line width=.5pt,line cap=round](0,0)--(\the\dimen@,0);}}$\hss}}}
\begingroup\lccode`\~=`\- \lowercase{\endgroup\let~}\lycee@moins
\newcommand\lycee@moinsepais{\mathcode`\-="8000 }
\AtBeginDocument{\mathchardef\lycee@moinsorig=\mathcode`\-\relax}
\AtBeginDocument{\catcode`\_=12 \mathcode`\_="8000
  \catcode`\^=12 \mathcode`\^="8000 }
\newcommand\lycee@exposants{%
  \fontdimen13\textfont2=.48\fontdimen6\textfont2
  \fontdimen14\textfont2=.48\fontdimen6\textfont2
  \fontdimen15\textfont2=.48\fontdimen6\textfont2
  \fontdimen13\scriptfont2=.48\fontdimen6\scriptfont2
  \fontdimen14\scriptfont2=.48\fontdimen6\scriptfont2
  \fontdimen15\scriptfont2=.48\fontdimen6\scriptfont2}
\every@math@size\expandafter{\the\every@math@size\lycee@exposants}
\newlength{\retraitcalculs}
\setlength{\retraitcalculs}{2em}

\RequirePackage{tikz}
\usetikzlibrary{arrows.meta,calc,babel,shapes.misc}
\RequirePackage[most]{tcolorbox}
\tcbuselibrary{skins,breakable}

\definecolor{coulDef}{HTML}{1F4E79}
\definecolor{coulProp}{HTML}{7030A0}
\definecolor{coulRem}{HTML}{C55A11}
\definecolor{coulEx}{HTML}{2E7D32}
\definecolor{coulExo}{HTML}{B03A2E}
\definecolor{coulPart}{HTML}{1F4E79}
\definecolor{coulMeth}{HTML}{1B6E4F}
\definecolor{coulCourbe}{HTML}{0B5ED7}
\definecolor{coulCourbeB}{HTML}{C00000}
\definecolor{coulCourbeC}{HTML}{2E7D32}
\definecolor{coulGrille}{HTML}{8C8C8C}
\definecolor{coulRep}{HTML}{1B6E4F}
\definecolor{coulBar}{HTML}{2E7D32}

\newcommand{\R}{\mathbb{R}}
\newcommand{\Rs}{\mathbb{R}^{*}}
\renewcommand{\le}{\leqslant}
\renewcommand{\ge}{\geqslant}
\renewcommand{\approx}{\simeq}
\newcommand{\vect}[1]{\overrightarrow{#1}}
\newcommand{\norme}[1]{\left\lVert #1 \right\rVert}
\newcommand{\intervalle}[2]{\left[#1\,;#2\right]}
\RequirePackage{cancel}
\renewcommand{\CancelColor}{\color{coulExo}}
\newcommand{\fois}[1]{\times{\color{coulExo}#1}}
\newcommand{\barre}[1]{\cancel{#1}}

\tikzset{grille/.style={coulGrille,line width=0.4pt}}
\newcommand{\quadrillage}[4]{\draw[grille,step=1] (#1,#2) grid (#3,#4);}
\tikzset{
  croix/.style={cross out,draw,line width=0.6pt,inner sep=0pt,minimum size=4pt},
  croixdroite/.style={croix,rotate=45,line width=0.8pt},
}
\newcommand{\pt}[3]{\path (#1) node[croix]{} node[#2,font=\small,inner sep=2.5pt]{$#3$};}

\newcommand{\lycee@repere}[4]{%
  \draw[grille,step=1] (#1,#2) grid (#3,#4);
  \draw[-{Stealth[length=2mm]},thick] (#1,0)--({#3+0.4},0) node[below left=-1pt]{$x$};
  \draw[-{Stealth[length=2mm]},thick] (0,#2)--(0,{#4+0.4}) node[below left=-1pt]{$y$};
  \node[below left,font=\scriptsize,inner sep=1.5pt] at (0,0) {$0$};}
\newcommand{\lycee@gradx}[1]{\foreach \k/\l in {#1}{%
  \draw (\k,0.07)--(\k,-0.07) node[below,font=\scriptsize,inner sep=1.5pt]{$\l$};}}
\newcommand{\lycee@grady}[1]{\foreach \k/\l in {#1}{%
  \draw (0.07,\k)--(-0.07,\k) node[left,font=\scriptsize,inner sep=1.5pt]{$\l$};}}
\newcommand{\lycee@intff}[2]{\left[#1\,;#2\right]}
\newcommand{\lycee@intfo}[2]{\left[#1\,;#2\right[}
\newcommand{\lycee@intof}[2]{\left]#1\,;#2\right]}
\newcommand{\lycee@intoo}[2]{\left]#1\,;#2\right[}
\newcommand{\lycee@ens}[1]{\left\{#1\right\}}
\AtBeginDocument{%
  \providecommand{\repere}{\lycee@repere}%
  \providecommand{\gradx}{\lycee@gradx}%
  \providecommand{\grady}{\lycee@grady}%
  \providecommand{\Cf}{\mathcal{C}\sb{\scriptscriptstyle f}}%
  \providecommand{\Cg}{\mathcal{C}\sb{\scriptscriptstyle g}}%
  \providecommand{\intff}{\lycee@intff}%
  \providecommand{\intfo}{\lycee@intfo}%
  \providecommand{\intof}{\lycee@intof}%
  \providecommand{\intoo}{\lycee@intoo}%
  \providecommand{\ens}{\lycee@ens}}

\newlist{questions}{enumerate}{3}
\setlist[questions,1]{label=\arabic*\textdegree),leftmargin=*,itemsep=2pt,topsep=3pt}
\setlist[questions,2]{label=\alph*),leftmargin=*,itemsep=1pt,topsep=2pt}
\setlist[questions,3]{label=\roman*),leftmargin=*,itemsep=1pt,topsep=1pt}
\setlist[itemize]{label=\textbullet,leftmargin=*,itemsep=2pt,topsep=3pt}

\newcommand{\besoinplace}[1]{\par\penalty\@M\begingroup
  \dimen@=#1\relax\dimen@ii=\pagegoal\advance\dimen@ii-\pagetotal
  \ifdim\dimen@>\dimen@ii\ifdim\dimen@ii>\z@\vfil\fi\break\fi\endgroup}

\newcommand{\lycee@pied}{\thepage/\pageref{LastPage}}
\newcommand{\entete}[2]{%
  \pagestyle{fancy}\fancyhf{}%
  \fancyhead[L]{\footnotesize\color{coulPart}\textbf{#1}}%
  \fancyhead[R]{\footnotesize\color{coulPart}\textbf{#2}}%
  \fancyfoot[C]{\footnotesize\lycee@pied}%
  \renewcommand{\headrulewidth}{0.4pt}%
  \renewcommand{\headrule}{\hbox to\headwidth{%
    \color{coulPart!40}\leaders\hrule height \headrulewidth\hfill}}}

\RequirePackage{ccicons}
\newcommand{\lycee@licence}{{\scriptsize\color{gray}%
  \copyright\ Luc Destombe --- \ccbyncsa\ CC BY-NC-SA 4.0}}
\newif\iflycee@licence \lycee@licencetrue
\newcommand{\sanslicence}{\lycee@licencefalse}
\AtBeginDocument{\iflycee@licence\AddToHookNext{shipout/foreground}{%
  \put(0,0){\rlap{\kern\dimexpr1in+\hoffset+\oddsidemargin+\textwidth\relax
    \llap{\raisebox{-\dimexpr1in+\voffset+\topmargin+\headheight+\headsep
      +\textheight+\footskip\relax}{\lycee@licence}}}}}\fi}

\newcommand{\formulesserrees}{%
  \setlength{\abovedisplayskip}{3pt}\setlength{\belowdisplayskip}{3pt}%
  \setlength{\abovedisplayshortskip}{2pt}\setlength{\belowdisplayshortskip}{3pt}}

\setlength{\parindent}{0pt}

\RequirePackage[implicit=false,hidelinks,pdfencoding=auto,psdextra,bookmarksopen,
  bookmarksopenlevel=1,pdfcreator={},pdfproducer={}]{hyperref}
\RequirePackage{bookmark}
\hypersetup{pdfauthor={Luc Destombe},pdfkeywords={CC BY-NC-SA 4.0}}
\newcounter{lycee@signet}
\newcommand{\signet}[3][]{\stepcounter{lycee@signet}%
  \edef\lycee@dest{\if\relax\detokenize{#1}\relax
    signet.\arabic{lycee@signet}\else#1\fi}%
  \hypertarget{\lycee@dest}{}\bookmark[level=#2,dest=\lycee@dest]{#3}}
\pdfstringdefDisableCommands{%
  \def\R{R}\def\vect#1{#1}\def\Cf{Cf}\def\Cg{Cg}\def\fois#1{×#1}%
  \def\barre#1{#1}\def\intervalle#1#2{[#1;#2]}\def\intff#1#2{[#1;#2]}%
  \def\textsuperscript#1{#1}\def\dfrac#1#2{#1/#2}\def\frac#1#2{#1/#2}%
  \def\vec#1{#1⃗}}%
\RequirePackage[official]{eurosym}

\geometry{top=0.8cm,bottom=1.3cm,left=1cm,right=1cm,footskip=0.6cm}

\pagestyle{fancy}
\fancyhf{}
\renewcommand{\headrulewidth}{0pt}
\fancyfoot[C]{\footnotesize Page \thepage{} / \pageref{LastPage}}

\newcommand{\titrefiche}[2]{%
  \begin{center}
    \Large\textbf{#1}\\[2pt]
    \large\textbf{#2}\\[2pt]
  \end{center}
  \vspace{0.10cm}}

\setlength{\columnseprule}{0.4pt}
\setlength{\columnsep}{15pt}
\renewcommand{\columnseprulecolor}{\color{black!45}}
\setlength{\multicolsep}{2pt}

\newcounter{partiefiche}
\renewcommand{\thepartiefiche}{\Roman{partiefiche}}
\newif\ifapresbandeau
\newcommand{\lycee@bandeau}[1]{%
  \par\penalty-600
  \vspace{0.08cm}%
  \noindent\colorbox{coulPart}{%
    \parbox{\dimexpr\linewidth-2\fboxsep\relax}{%
      \color{white}\bfseries\raggedright\strut#1\strut}}%
  \nopagebreak\par\nopagebreak
  \vspace{0.06cm}\nopagebreak
  \apresbandeautrue}
\newcommand{\partiefiche}[1]{%
  \refstepcounter{partiefiche}%
  \lycee@bandeau{\signet{1}{\thepartiefiche{} --- #1}\thepartiefiche{} --- #1}}
\newcommand{\partiefichesansnum}[1]{\lycee@bandeau{\signet{1}{#1}#1}}

\newcounter{exo}
\newlength{\espaceexo}\setlength{\espaceexo}{7pt}
\newlength{\espacetitre}\setlength{\espacetitre}{2pt}
\newcommand{\exo}[1][\relax]{%
  \par
  \ifapresbandeau\apresbandeaufalse\else\penalty-150\addvspace{\espaceexo}\fi
  \refstepcounter{exo}%
  \noindent\signet[exercice-\arabic{exo}]{2}{Exercice \arabic{exo}}%
  \textbf{\color{coulExo}\underline{Exercice \theexo}}%
  \ifx\relax#1\relax\else\textbf{ : #1}\fi
  \par\nopagebreak\vspace{\espacetitre}\nopagebreak
  \ignorespaces}

\setlist[enumerate]{nosep,leftmargin=*,itemsep=2pt,topsep=2pt}
\setlist[itemize]{nosep,leftmargin=*,topsep=2pt}
\newcommand{\choix}[3]{\par\hspace*{1em}%
  \makebox[0.3\linewidth][l]{a) #1}\makebox[0.3\linewidth][l]{b) #2}c) #3}
\newcommand{\deux}[2]{\par\hspace*{0.6em}%
  \makebox[0.48\linewidth][l]{#1}#2}
\newcommand{\trois}[3]{\par\hspace*{0.6em}%
  \makebox[0.32\linewidth][l]{#1}\makebox[0.32\linewidth][l]{#2}#3}

  \renewcommand{\exo}[1][\relax]{%
    \ifapresbandeau\apresbandeaufalse\else\par\penalty-150\fi
    \refstepcounter{exo}%
    \vspace{0.05cm}%
    \noindent\signet[exercice-\arabic{exo}]{2}{Exercice \arabic{exo}}%
    \textbf{\color{coulExo}\underline{Exercice \theexo}}%
    \ifx\relax#1\relax\else\textbf{ : #1}\fi%
    \par\nopagebreak
    \ignorespaces}
  \newcommand{\rappel}[1]{%
    \par\vspace{0.06cm}\nopagebreak
    \noindent\fcolorbox{black!35}{black!5}{%
      \parbox{\dimexpr\linewidth-2\fboxsep-2\fboxrule\relax}{%
        \small\textbf{Rappel.} #1}}%
    \par\vspace{0.06cm}\nopagebreak}
  \setlist[enumerate]{nosep,leftmargin=*}
  \setlist[itemize]{nosep,leftmargin=*}
  \setlength{\multicolsep}{3pt plus 1pt minus 1pt}
  \setlength{\premulticols}{10pt}
  \setlength{\postmulticols}{10pt}
  \newlist{expr}{enumerate}{1}
  \setlist[expr]{label={\Alph*\ $=$},nosep,leftmargin=*,itemsep=0pt}

\renewenvironment{center}{\par\addvspace{3pt}\centering}{\par\addvspace{3pt}}
\formulesserrees
\AtBeginDocument{\formulesserrees}
\clubpenalty=10000
\widowpenalty=10000
\displaywidowpenalty=10000
\brokenpenalty=10000
\setlength{\parskip}{0pt}

\RequirePackage{ragged2e}
\setlength{\RaggedRightRightskip}{0pt plus 1fil}
\AtBeginDocument{\RaggedRight\binoppenalty=10000 \relpenalty=10000 }
\g@addto@macro\@parboxrestore{\RaggedRight}
\makeatother
\usepackage{tkz-tab}

\begin{document}

\titrefiche{1\textsuperscript{re} STI2D --- Dérivation et étude de fonctions}{Fiche d'exercices du chapitre}

\begin{multicols*}{2}\raggedcolumns

\partiefiche{Fonctions affines et droites}

\exo[Coefficient directeur]
Pour chaque couple de points, calculer le coefficient directeur de la droite $(AB)$.
\begin{enumerate}[label=\alph*)]
  \item $A(1\,;2)$ et $B(3\,;8)$
  \item $A(-1\,;4)$ et $B(2\,;-2)$
  \item $A(0\,;-3)$ et $B(4\,;-1)$
  \item $A(-3\,;5)$ et $B(1\,;5)$
\end{enumerate}

\exo[Équation d'une droite]
Pour chaque couple de points, déterminer l'équation réduite de la droite $(AB)$.
\begin{enumerate}[label=\alph*)]
  \item $A(1\,;1)$ et $B(2\,;4)$
  \item $A(-1\,;5)$ et $B(1\,;1)$
  \item $A(-2\,;-3)$ et $B(2\,;5)$
  \item $A(2\,;0)$ et $B(6\,;2)$
\end{enumerate}

\exo[Une fonction affine par deux images]
\begin{enumerate}[label=\arabic*)]
  \item $f$ est affine, avec $f(2)=7$ et $f(-1)=-2$. Déterminer $f(x)$, puis calculer
  $f(10)$.
  \item $g$ est affine, avec $g(-3)=8$ et $g(1)=0$. Déterminer $g(x)$, puis calculer
  $g(5)$.
\end{enumerate}

\exo[Lire des équations]
Lire l'équation réduite de chacune des droites $(d_{1})$, $(d_{2})$ et $(d_{3})$.
\begin{center}
\begin{tikzpicture}[x=0.45cm,y=0.45cm]
  \repere{-3}{-3}{5}{4}
  \gradx{-2,-1,1,2,3,4} \grady{-2,-1,1,2,3}
  \begin{scope}
  \clip (-3,-3) rectangle (5,4);
  \draw[coulCourbe,line width=1.1pt] (-1,-3)--(2.5,4);
  \draw[coulCourbeB,line width=1.1pt] (-1,4)--(5,-2);
  \draw[coulCourbeC,line width=1.1pt] (-2,-3)--(5,0.5);
  \end{scope}
  \node[coulCourbe,font=\footnotesize] at (1.4,3.6) {$(d_{1})$};
  \node[coulCourbeB,font=\footnotesize] at (-0.6,2.7) {$(d_{2})$};
  \node[coulCourbeC,font=\footnotesize] at (4.2,0.8) {$(d_{3})$};
\end{tikzpicture}
\end{center}

\partiefiche{Nombre dérivé}

\exo[Nombre dérivé en $1$]
Soit $f(x)=x^{2}+2x-1$.
\begin{enumerate}[label=\alph*)]
  \item Calculer $f(1)$, puis $f(1+h)$.
  \item Simplifier le taux de variation $T(h)=\dfrac{f(1+h)-f(1)}{h}$.
  \item En déduire $f'(1)$.
\end{enumerate}

\exo[Nombre dérivé en $2$]
Soit $g(x)=2x^{2}-3x$. Calculer $g(2)$ et $g(2+h)$, simplifier
le taux de variation $T(h)=\dfrac{g(2+h)-g(2)}{h}$, puis en déduire $g'(2)$.

\exo[Avec un coefficient négatif]
Soit $k(x)=-x^{2}+5x$. Calculer $k(1)$ et $k(1+h)$, simplifier
le taux de variation $T(h)=\dfrac{k(1+h)-k(1)}{h}$, puis en déduire $k'(1)$.

\exo[À la calculatrice]
Soit $f(x)=x^{2}$ et $a=4$.
\begin{enumerate}[label=\alph*)]
  \item Compléter le tableau, où $T(h)=\dfrac{f(4+h)-f(4)}{h}$.
  \begin{center}
  \renewcommand{\arraystretch}{1.4}
  \begin{tabular}{|c|c|c|c|}
  \hline
  $h$ & $1$ & $0{,}1$ & $0{,}01$\\
  \hline
  $T(h)$ & & & \\
  \hline
  \end{tabular}
  \end{center}
  \item Quelle valeur de $f'(4)$ ce tableau suggère-t-il ?
\end{enumerate}

\partiefiche{Tangente à une courbe}

\exo[Lire des nombres dérivés]
\begin{minipage}[c]{0.42\linewidth}\raggedright
On a tracé $\Cf$ et ses tangentes aux points $A$, $B$, $C$ et $E$.

\smallskip
Lire $f'(-3)$, $f'(-1)$, $f'(1)$ et $f'(4)$.
\end{minipage}\hfill
\begin{minipage}[c]{0.56\linewidth}\centering
\begin{tikzpicture}[x=0.45cm,y=0.45cm]
  \repere{-4}{-2}{5}{4}
  \gradx{-3,-2,-1,1,2,3,4} \grady{-1,1,2,3}
  \draw[coulCourbe,line width=1.1pt]
    (-3,2) .. controls (-2.333,0.667) and (-1.667,-1) .. (-1,-1)
           .. controls (-0.333,-1) and (0.333,-0.333) .. (1,1)
           .. controls (1.667,2.333) and (2.333,3) .. (3,3)
           .. controls (3.333,3) and (3.667,2.667) .. (4,2);
  \draw[coulCourbeB,thick] (-3.8,3.6)--(-2.2,0.4);
  \draw[coulCourbeB,thick] (-2,-1)--(0,-1);
  \draw[coulCourbeB,thick] (0.2,-0.6)--(1.8,2.6);
  \draw[coulCourbeB,thick] (3.4,3.2)--(4.6,0.8);
  \pt{-3,2}{left}{A}
  \pt{-1,-1}{below}{B}
  \pt{1,1}{right}{C}
  \pt{4,2}{right}{E}
  \node[coulCourbe,font=\footnotesize] at (2.2,3.5) {$\Cf$};
\end{tikzpicture}
\end{minipage}

\exo[Tracer des tangentes]
Dans un repère, placer chaque point de $\Cf$, puis tracer la tangente à $\Cf$ en ce
point.
\begin{enumerate}[label=\alph*)]
  \item $A(0\,;1)$ et $f'(0)=2$
  \item $B(1\,;3)$ et $f'(1)=-1$
  \item $C(-2\,;-1)$ et $f'(-2)=\frac{1}{2}$
  \item $D(3\,;-2)$ et $f'(3)=0$
\end{enumerate}

\exo[Retrouver l'allure d'une courbe]
Dans un repère, tracer une courbe qui passe par $A(-3\,;-2)$, $B(-1\,;2)$,
$C(1\,;0)$ et $D(3\,;-2)$, avec $f'(-3)=3$, $f'(-1)=0$, $f'(1)=-2$ et $f'(3)=0$.
Commencer par tracer les quatre tangentes.

\partiefiche{Équation de la tangente}

\exo[Avec un nombre dérivé admis]
Soit $f(x)=x^{2}-3x+2$. On admet que $f'(4)=5$ et $f'(1)=-1$.
\begin{enumerate}[label=\alph*)]
  \item Calculer $f(4)$, puis déterminer l'équation réduite de la tangente à $\Cf$
  au point d'abscisse $4$.
  \item Même question au point d'abscisse $1$.
\end{enumerate}

\exo[Quatre tangentes]
Pour chaque cas, déterminer l'équation réduite de la tangente à $\Cf$ au point
d'abscisse $a$.
\begin{enumerate}[label=\alph*)]
  \item $a=2$ ; $f(2)=3$ et $f'(2)=1$
  \item $a=-1$ ; $f(-1)=4$ et $f'(-1)=-2$
  \item $a=3$ ; $f(3)=-1$ et $f'(3)=0$
  \item $a=0$ ; $f(0)=5$ et $f'(0)=-3$
\end{enumerate}

\exo[Du taux de variation à la tangente]
Soit $f(x)=-x^{2}+3x+1$.
\begin{enumerate}[label=\alph*)]
  \item Calculer $f(2)$, puis $f(2+h)$.
  \item Simplifier le taux de variation $T(h)=\dfrac{f(2+h)-f(2)}{h}$ et en déduire
  $f'(2)$.
  \item Déterminer l'équation réduite de la tangente à $\Cf$ au point d'abscisse $2$.
\end{enumerate}

\partiefiche{Approximation affine}

\exo[Le carré de $3{,}02$]
Soit $f(x)=x^{2}$. Sa tangente au point d'abscisse $3$ a pour équation $y=6x-9$.
\begin{enumerate}[label=\alph*)]
  \item Avec la tangente, donner une valeur approchée de $f(3{,}02)$, puis calculer
  $f(3{,}02)$ et comparer.
  \item Même question pour $f(2{,}99)$.
\end{enumerate}

\exo[Une fonction inverse]
Soit $f(x)=\dfrac{1}{x-1}$ pour $x\neq 1$. On admet que $f'(2)=-1$.
\begin{enumerate}[label=\alph*)]
  \item Calculer $f(2)$, puis déterminer l'équation de la tangente au point
  d'abscisse $2$.
  \item Avec la tangente, donner une valeur approchée de $f(2{,}1)$.
  \item Calculer $f(2{,}1)$ à la calculatrice et comparer.
\end{enumerate}

\exo[Une plaque qui se dilate]
Une plaque carrée de côté $c$ (en cm) a pour aire $A(c)=c^{2}$. On admet que
$A'(10)=20$. Chauffée, la plaque passe de $10$~cm à $10{,}1$~cm de côté.
\begin{enumerate}[label=\alph*)]
  \item Déterminer l'équation de la tangente à la courbe de $A$ au point
  d'abscisse $10$.
  \item En déduire une valeur approchée de l'aire de la plaque chauffée.
  \item Calculer l'aire exacte et comparer.
\end{enumerate}

\partiefiche{Fonction dérivée d'une fonction polynôme}

\exo[Dériver]
Calculer la dérivée de chaque fonction.
\begin{multicols}{2}
\begin{enumerate}[label=\alph*)]
  \item $f(x)=7x-5$
  \item $g(x)=1-4x$
  \item $h(x)=x^{2}+6x-2$
  \item $k(x)=3x^{2}-x+8$
  \item $m(x)=-5x^{2}+10x$
  \item $p(x)=x^{3}+4x^{2}-x$
  \item $q(x)=2x^{4}-3x^{2}+1$
  \item $r(x)=-x^{5}+2x^{3}$
\end{enumerate}
\end{multicols}

\exo[Degré 3 et plus]
Calculer la dérivée de chaque fonction.
\begin{enumerate}[label=\alph*)]
  \item $f(x)=4x^{3}-6x^{2}+5x-1$
  \item $g(x)=x^{6}-x^{4}$
  \item $h(x)=-2x^{5}+x^{3}-7x$
  \item $k(x)=5-3x^{4}$
  \item $m(x)=\frac{1}{3}x^{3}-2x$
\end{enumerate}

\exo[Entraînement]
Calculer la dérivée de chaque fonction.
\begin{multicols}{2}
\begin{enumerate}[label=\alph*)]
  \item $f(x)=6x-2$
  \item $g(x)=-x^{2}+4x-3$
  \item $h(x)=5x^{3}-2x^{2}$
  \item $k(x)=x^{4}+x^{3}+x^{2}+x$
  \item $m(x)=-3x^{4}+6x$
  \item $p(x)=2x^{5}-5x^{2}+3$
  \item $q(x)=\frac{1}{2}x^{2}+3x$
  \item $r(x)=\frac{1}{4}x^{4}-x$
\end{enumerate}
\end{multicols}

\exo[Développer d'abord]
Développer chaque expression, puis calculer sa dérivée.
\begin{multicols}{2}
\begin{enumerate}[label=\alph*)]
  \item $A(x)=(x-1)(x+4)$
  \item $B(x)=(2x+1)(x-3)$
  \item $C(x)=(x+2)^{2}$
  \item $D(x)=x(x^{2}-5)$
\end{enumerate}
\end{multicols}

\exo[Dérivée et tangentes]
Soit $f(x)=x^{3}-2x+1$.
\begin{enumerate}[label=\alph*)]
  \item Calculer $f'(x)$.
  \item Calculer $f(1)$ et $f'(1)$, puis déterminer l'équation de la tangente au point
  d'abscisse $1$.
  \item Même question au point d'abscisse $0$.
\end{enumerate}

\partiefiche{Autres fonctions dérivées}

\exo[Inverse, cosinus, sinus]
Calculer la dérivée de chaque fonction.
\begin{multicols}{2}
\begin{enumerate}[label=\alph*)]
  \item $f(x)=\dfrac{4}{x}$
  \item $g(x)=-\dfrac{7}{x}$
  \item $h(x)=\cos x+3x$
  \item $k(x)=5\sin x$
  \item $m(x)=2\cos x-3\sin x$
  \item $p(x)=x^{4}-\dfrac{1}{x}$
  \item $q(x)=\dfrac{6}{x}+2x^{2}$
  \item $r(x)=-\sin x+x^{3}$
\end{enumerate}
\end{multicols}

\exo[Entraînement]
Calculer la dérivée de chaque fonction.
\begin{multicols}{2}
\begin{enumerate}[label=\alph*)]
  \item $f(x)=\dfrac{3}{x}+x$
  \item $g(x)=4\cos x+x^{2}$
  \item $h(x)=-2\sin x+5$
  \item $k(x)=x^{3}-\dfrac{5}{x}$
  \item $m(x)=\sin x-\cos x$
  \item $p(x)=\dfrac{1}{x}-3\cos x$
\end{enumerate}
\end{multicols}

\exo[Tangentes à la courbe de $\frac{2}{x}$]
Soit $f(x)=\dfrac{2}{x}$ pour $x\neq 0$.
\begin{enumerate}[label=\alph*)]
  \item Calculer $f'(x)$.
  \item Déterminer l'équation de la tangente au point d'abscisse $1$.
  \item Même question au point d'abscisse $2$.
\end{enumerate}

\exo[Sinus et cosinus près de $0$]
\begin{enumerate}[label=\alph*)]
  \item Soit $f(x)=\sin x$. Calculer $f'(0)$, puis déterminer l'équation de la
  tangente au point d'abscisse $0$.
  \item En déduire une valeur approchée de $\sin(0{,}1)$ ; comparer avec la
  calculatrice (en radians).
  \item Soit $g(x)=\cos x$. Calculer $g'(0)$. Que peut-on dire de la tangente au point
  d'abscisse $0$ ?
\end{enumerate}

\partiefiche{Dérivée d'un produit, d'un carré}

\exo[Des carrés]
Calculer la dérivée de chaque fonction, avec la formule de $u^{2}$.
\begin{multicols}{2}
\begin{enumerate}[label=\alph*)]
  \item $f(x)=(4x+1)^{2}$
  \item $g(x)=(2-5x)^{2}$
  \item $h(x)=(x^{2}+3)^{2}$
  \item $k(x)=(x^{3}-1)^{2}$
\end{enumerate}
\end{multicols}

\exo[Des produits]
Calculer la dérivée de chaque fonction, avec la formule de $u\times v$.
\begin{enumerate}[label=\alph*)]
  \item $f(x)=(x+2)(3x-1)$
  \item $g(x)=(2x-3)(x^{2}+1)$
  \item $h(x)=x^{2}(4x+5)$
  \item $k(x)=(x^{2}-x)(x^{2}+2)$
  \item $m(x)=x\sin x$
\end{enumerate}

\exo[Entraînement]
Calculer la dérivée de chaque fonction.
\begin{multicols}{2}
\begin{enumerate}[label=\alph*)]
  \item $f(x)=(3x+1)(x-2)$
  \item $g(x)=(x+5)^{2}$
  \item $h(x)=(x^{2}-4)(2x+3)$
  \item $k(x)=(3-x^{2})^{2}$
  \item $m(x)=x^{2}\cos x$
  \item $p(x)=(2x-1)(x^{3}+x)$
\end{enumerate}
\end{multicols}

\exo[Deux façons de faire]
Soit $f(x)=(x+3)(2x-1)$.
\begin{enumerate}[label=\alph*)]
  \item Calculer $f'(x)$ avec la formule du produit.
  \item Développer $f(x)$, puis le dériver.
  \item Comparer les deux résultats.
\end{enumerate}

\partiefiche{Dérivée de l'inverse, d'un quotient}

\exo[Des inverses]
Calculer la dérivée de chaque fonction.
\begin{multicols}{2}
\begin{enumerate}[label=\alph*)]
  \item $f(x)=\dfrac{1}{3x+2}$
  \item $g(x)=\dfrac{1}{1-4x}$
  \item $h(x)=\dfrac{1}{x^{2}+1}$
  \item $k(x)=\dfrac{5}{2x-1}$
\end{enumerate}
\end{multicols}

\exo[Des quotients]
Calculer la dérivée de chaque fonction.
\begin{multicols}{2}
\begin{enumerate}[label=\alph*)]
  \item $f(x)=\dfrac{x+4}{x-1}$
  \item $g(x)=\dfrac{2x-1}{x+2}$
  \item $h(x)=\dfrac{3x}{x+1}$
  \item $k(x)=\dfrac{x^{2}+1}{x-2}$
\end{enumerate}
\end{multicols}

\exo[Entraînement]
Calculer la dérivée de chaque fonction.
\begin{multicols}{2}
\begin{enumerate}[label=\alph*)]
  \item $f(x)=\dfrac{1}{x+5}$
  \item $g(x)=\dfrac{3}{4x-1}$
  \item $h(x)=\dfrac{x-3}{x+2}$
  \item $k(x)=\dfrac{4x+1}{2x-3}$
  \item $m(x)=\dfrac{2}{x^{2}+4}$
  \item $p(x)=\dfrac{x^{2}}{x+1}$
\end{enumerate}
\end{multicols}

\exo[Quotient et tangente]
Soit $f(x)=\dfrac{2x+1}{x-1}$ pour $x\neq 1$.
\begin{enumerate}[label=\alph*)]
  \item Calculer $f'(x)$.
  \item Calculer $f(2)$ et $f'(2)$, puis déterminer l'équation de la tangente au point
  d'abscisse $2$.
\end{enumerate}

\partiefiche{Fonctions composées}

\exo[Des puissances]
Calculer la dérivée de chaque fonction, avec la formule de $u^{n}$.
\begin{multicols}{2}
\begin{enumerate}[label=\alph*)]
  \item $f(x)=(3x-2)^{4}$
  \item $g(x)=(x^{2}+1)^{3}$
  \item $h(x)=(1-2x)^{5}$
  \item $k(x)=(x^{3}+x)^{2}$
\end{enumerate}
\end{multicols}

\exo[Cosinus et sinus composés]
Calculer la dérivée de chaque fonction.
\begin{multicols}{2}
\begin{enumerate}[label=\alph*)]
  \item $f(x)=\cos(4x)$
  \item $g(x)=\sin(2x+3)$
  \item $h(x)=3\cos(5x-1)$
  \item $k(t)=2\sin(100t)$
\end{enumerate}
\end{multicols}

\exo[Entraînement]
Calculer la dérivée de chaque fonction.
\begin{multicols}{2}
\begin{enumerate}[label=\alph*)]
  \item $f(x)=(5x+2)^{3}$
  \item $g(x)=(x^{2}-3x)^{4}$
  \item $h(x)=\cos(3x-2)$
  \item $k(x)=4\sin(x^{2})$
  \item $m(t)=5\cos(100t+1)$
  \item $p(x)=(2-x)^{6}$
\end{enumerate}
\end{multicols}

\exo[Un piston]
La position d'un piston, en cm, au bout de $t$ secondes, est $p(t)=6\cos(2t)$. Sa
vitesse, en cm/s, est la dérivée $p'(t)$.
\begin{enumerate}[label=\alph*)]
  \item Calculer $p'(t)$.
  \item Quelle est la vitesse du piston à l'instant $t=0$ ?
  \item Sachant que $\sin(2t)$ est toujours compris entre $-1$ et $1$, quelle est la
  plus grande vitesse possible du piston (en valeur absolue) ?
\end{enumerate}

\exo[Dérivées en vrac]
Pour chaque fonction, choisir la bonne formule, puis calculer la dérivée.
\begin{multicols}{2}
\begin{enumerate}[label=\alph*)]
  \item $f(x)=3x^{4}-\dfrac{2}{x}$
  \item $g(x)=(x^{2}+1)(x-3)$
  \item $h(x)=(2x+5)^{3}$
  \item $k(x)=\dfrac{4}{x^{2}+3}$
  \item $m(x)=\dfrac{x-1}{x^{2}+1}$
  \item $p(x)=x\cos x$
\end{enumerate}
\end{multicols}

\partiefiche{Dérivée et sens de variation}

\exo[Du signe de $f'$ aux variations]
Voici le signe de $f'(x)$ sur $\intff{-4}{5}$. On sait que $f(-4)=2$, $f(-1)=-3$,
$f(3)=4$ et $f(5)=1$.
\begin{center}
\begin{tikzpicture}[scale=0.8]
  \tkzTabInit[lgt=1.6,espcl=1.8]{$x$/1,$f'(x)$/1}{$-4$,$-1$,$3$,$5$}
  \tkzTabLine{,-,z,+,z,-,}
\end{tikzpicture}
\end{center}
Dresser le tableau de variations de $f$.

\exo[Un trinôme]
Soit $f(x)=x^{2}-6x+5$ sur $\intff{0}{6}$.
\begin{enumerate}[label=\alph*)]
  \item Calculer $f'(x)$ et étudier son signe.
  \item Dresser le tableau de variations de $f$.
\end{enumerate}

\exo[Un polynôme de degré 3]
Soit $f(x)=x^{3}-3x^{2}+1$ sur $\intff{-1}{3}$.
\begin{enumerate}[label=\alph*)]
  \item Calculer $f'(x)$.
  \item Étudier le signe de $f'(x)$.
  \item Dresser le tableau de variations de $f$.
\end{enumerate}

\exo[Un coefficient négatif]
Soit $g(x)=-2x^{3}+3x^{2}+12x-5$ sur $\intff{-2}{3}$.
\begin{enumerate}[label=\alph*)]
  \item Calculer $g'(x)$.
  \item Étudier le signe de $g'(x)$.
  \item Dresser le tableau de variations de $g$.
\end{enumerate}

\exo[Un quotient]
Soit $f(x)=\dfrac{2x+1}{x+1}$ sur $\intff{0}{4}$.
\begin{enumerate}[label=\alph*)]
  \item Montrer que $f'(x)=\dfrac{1}{(x+1)^{2}}$.
  \item En déduire le sens de variation de $f$ sur $\intff{0}{4}$, puis dresser son
  tableau de variations.
\end{enumerate}

\exo[Étude complète et tracé]
Soit $f(x)=x^{3}-3x$ sur $\intff{-2}{2}$.
\begin{enumerate}[label=\alph*)]
  \item Calculer $f'(x)$ et étudier son signe.
  \item Dresser le tableau de variations de $f$.
  \item Déterminer l'équation de la tangente $T$ au point d'abscisse $0$.
  \item Compléter un tableau de valeurs de $f$ pour $x=-2$, $-1$, $0$, $1$ et $2$.
  \item Dans un repère, tracer $T$, puis $\Cf$.
\end{enumerate}

\partiefiche{Problèmes concrets : maximum et minimum}

\exo[Un coût minimal]
Une usine fabrique $q$ dizaines de pièces par jour, avec $0\le q\le 25$. Le coût de
fabrication, en euros, est $C(q)=q^{2}-30q+300$.
\begin{enumerate}[label=\alph*)]
  \item Calculer $C'(q)$.
  \item Dresser le tableau de variations de $C$ sur $\intff{0}{25}$.
  \item Combien de pièces faut-il fabriquer pour que le coût soit minimal ? Que
  vaut-il ?
\end{enumerate}

\exo[La boîte la plus grande]
Dans un carton carré de $12$~cm de côté, on découpe aux quatre coins un carré de côté
$x$ (en cm), puis on relève les bords pour former une boîte sans couvercle. Son volume
est $V(x)=x(12-2x)^{2}$, pour $0\le x\le 6$.
\begin{enumerate}[label=\alph*)]
  \item Montrer que $V(x)=4x^{3}-48x^{2}+144x$.
  \item Calculer $V'(x)$.
  \item Étudier le signe de $V'(x)$ sur $\intff{0}{6}$.
  \item Dresser le tableau de variations de $V$.
  \item Pour quelle valeur de $x$ le volume est-il maximal ? Que vaut-il ?
\end{enumerate}

\exo[Le tambour d'une machine à laver]
Après l'essorage, le tambour d'une machine à laver tourne librement jusqu'à l'arrêt.
Au bout de $t$ secondes, il a fait $N(t)=-0{,}5t^{2}+20t$ tours, pour
$0\le t\le 40$. Sa vitesse de rotation, en tours par seconde, est $N'(t)$.
\begin{enumerate}[label=\alph*)]
  \item Calculer $N'(t)$.
  \item Au bout de combien de secondes le tambour s'arrête-t-il ?
  \item Combien de tours a-t-il faits jusqu'à l'arrêt ?
\end{enumerate}

\partiefiche{Résoudre une équation $f(x)=k$}

\exo[Nombre de solutions]
Voici le tableau de variations d'une fonction $f$ dérivable sur $\intff{-3}{4}$.
\begin{center}
\begin{tikzpicture}[scale=0.8]
  \tkzTabInit[lgt=1.4,espcl=1.8]{$x$/1,$f(x)$/1.6}{$-3$,$0$,$2$,$4$}
  \tkzTabVar{-/$-2$,+/$5$,-/$-4$,+/$1$}
\end{tikzpicture}
\end{center}
Pour chaque équation, donner le nombre de solutions et un intervalle qui contient
chacune d'elles.
\begin{multicols}{2}
\begin{enumerate}[label=\alph*)]
  \item $f(x)=0$
  \item $f(x)=3$
  \item $f(x)=-3$
  \item $f(x)=6$
\end{enumerate}
\end{multicols}

\exo[Une seule solution]
Soit $f(x)=x^{3}+x-3$ sur $\intff{0}{2}$.
\begin{enumerate}[label=\alph*)]
  \item Montrer que $f$ est strictement croissante sur $\intff{0}{2}$.
  \item Montrer que l'équation $f(x)=0$ a une seule solution $\alpha$ dans
  $\intff{0}{2}$.
  \item Avec la calculatrice, encadrer $\alpha$ entre deux nombres à un chiffre après
  la virgule.
\end{enumerate}

\besoinplace{9\baselineskip}
\exo[Trois solutions]
Soit $g(x)=x^{3}-12x$ sur $\intff{-4}{4}$.
\begin{enumerate}[label=\alph*)]
  \item Calculer $g'(x)$ et étudier son signe.
  \item Dresser le tableau de variations de $g$.
  \item Combien de solutions l'équation $g(x)=5$ a-t-elle dans $\intff{-4}{4}$ ?
  \item Encadrer la plus grande d'entre elles entre deux nombres à un chiffre après la
  virgule.
\end{enumerate}

\end{multicols*}

\end{document}
